(1) 令介质相对实验室速度为 \(+u\)。相对论平行速度合成给顺流相速度 \(v_+=(c/n+u)/(1+u/(nc))\)，逆流光的有符号速度为 \(v_-=(-c/n+u)/(1-u/(nc))\)，在 \(u<c/n\) 的题设范围内后者为负。使用简单 Galilean 相加会漏掉分母中光速不变性修正。
(2) 展开 \((1+\epsilon)^{-1}=1-\epsilon+O(\epsilon^2)\)，得 \(v_+=c/n+u(1-1/n^2)+O(u^2/c)\)，\(v_-=-c/n+u(1-1/n^2)+O(u^2/c)\)，故顺流波在实验室中的一阶增速不是全额 \(u\)，而是 \(u(1-1/n^2)\)。
(3) 在实验室固定的有效传播长度 \(L\) 内，顺流时间 \(t_+=L/v_+\)，逆流时间 \(t_-=L/|v_-|\)。一阶倒数展开得到 \(\Delta t=t_--t_+=2Lu(n^2-1)/c^2+O(u^2)\)；若合束比较同频单色光的实验室相位，以真空波长 \(\lambda_0\) 记，其对应条纹差为 \(\Delta N=(c/\lambda_0)\Delta t\simeq2Lu(n^2-1)/(c\lambda_0)\)，这里忽略系统其他臂和端点的相位贡献。
(4) \(n\to1\) 时拖曳系数及时间差趋零，\(u\to0\) 时两向时间相等。若介质强色散，光相速度和脉冲群速度的拖曳公式不再等同，正确表达需要折射率随频率变换；实验中管端、分束与反射面是否随介质移动也会改变时延及观测相位，不能把固定端点假设隐去。
